%%% SECTION 3: CALIFORNIA AND FEDERAL REGULATORY LANDSCAPE %%%

[This section should be approximately 10-12 pages. The major theme is a
comprehensive dual-jurisdiction analysis of how California state law and
federal law together govern AI and robotics in oncology clinical trials.
California is the site of the first Physical AI oncology clinical trial, making
this dual analysis essential for understanding the complete regulatory
environment.]

[PRIMARY SOURCE: national-platform/research\_b/01\_research\_b\_part1.txt and
02\_research\_b\_part2.txt. These files contain the complete California and
federal legal framework analysis. All content should be adapted from research\_b
to Physical AI oncology trial contexts.]

[IMPORTANT: Use research\_b to fact-check all bill names referenced in the
Clinical Trial Site Complete Documentation \cite{site-docs}
(national-platform/new\_trial\_psl/ files). Verify that SB 1042, AB 2847,
SB 892, and other legislative references use correct bill names and governance
structures. If any discrepancies exist between the PSL documents and
research\_b's authoritative analysis, note corrections and update titles
and content accordingly.]

\subsection{Regulatory Classification Framework}
\label{subsec:regulatory-classification}

[Adapt the regulatory map from research\_b/01\_research\_b\_part1.txt. Cover the
two-axis classification system: (i) regulated medical product vs. research
operational tool, and (ii) investigational vs. approved and significant risk
vs. non-significant risk for devices. Explain how Physical AI systems may
simultaneously occupy multiple classification categories.]

[Reference the federal device definition (21 U.S.C. \S 321(h)), the device
investigational exemption authority (21 U.S.C. \S 360j(g)), and the FD\&C Act
\S 520(o) software exclusions from the 21st Century Cures Act as described in
research\_b.]

\subsection{California Regulatory Authority}
\label{subsec:california-authority}

[Adapt the California authority catalog from research\_b/01\_research\_b\_part1.txt.
Organize by regulatory domain:]

\subsubsection{Human Subjects Protections}
[Cover the Protection of Human Subjects in Medical Experimentation Act
(Cal. Health \& Safety Code \S\S 24170-24179.5), the Experimental Subject's
Bill of Rights (\S 24172), and the definition of "medical experiment"
(\S 24174). Explain how these provisions create additional protections beyond
federal requirements for Physical AI trial participants in California.]

\subsubsection{Health Data Privacy and Security}
[Cover the Confidentiality of Medical Information Act (CMIA, Cal. Civ. Code
\S\S 56-56.37), electronic medical information safeguards (\S 56.101),
California data breach notification (\S 1798.82), and medical information
breach requirements (Health \& Safety Code \S 1280.15). Explain how these
provisions interact with HIPAA \cite{hipaa} in Physical AI trial contexts
where robots generate and process health data.]

\subsubsection{Consumer Privacy}
[Cover CCPA/CPRA (Cal. Civ. Code \S\S 1798.100 et seq.) and its clinical
trial/biomedical research exemptions. Explain the practical implications for
Physical AI data processing, noting that clinical trial data may be statutorily
exempt in many common trial contexts but that non-exempt data streams from
robotic systems may still trigger CCPA obligations.]

\subsubsection{AI-Specific Healthcare Communications}
[Cover AB 3030 (Cal. Health \& Safety Code \S 1339.75) on AI-generated patient
clinical communications disclosure and AB 489 on health advice from artificial
intelligence. Explain how Physical AI systems that communicate with patients
must comply with these disclosure requirements.]

\subsubsection{Medical Practice and Licensing}
[Cover the unlicensed practice of medicine provisions (Cal. Bus. \& Prof. Code
\S 2052) and how they apply to Physical AI systems that may perform functions
traditionally reserved for licensed practitioners.]

\subsection{Federal Regulatory Authority}
\label{subsec:federal-authority}

[Adapt the federal authority catalog from research\_b/01\_research\_b\_part1.txt.
Organize by regulatory domain:]

\subsubsection{Investigational Product Frameworks}
[Cover IND regulations (21 CFR Part 312) \cite{cfr-part-312}, IDE regulations
(21 CFR Part 812), and the interaction between drug and device regulation when
Physical AI systems are used alongside investigational drugs. Reference the
Physical AI Adaption: 21 CFR Part 312 \cite{cfr312-adapt} as the adapted
framework.]

\subsubsection{Human Subjects Protection}
[Cover 21 CFR Part 50 \cite{cfr-part-50} (FDA human subjects protection),
21 CFR Part 56 (IRBs), and the Common Rule (45 CFR Part 46). Explain how
Physical AI adds new dimensions to informed consent, ongoing monitoring, and
IRB review. Reference Physical AI Adaption: 21 CFR Part 50 \cite{cfr50-adapt}
and the NCI CIRB \cite{nci-cirb}.]

\subsubsection{Electronic Records and Data Integrity}
[Cover 21 CFR Part 11 (electronic records and signatures) as it applies to
Physical AI systems that generate, store, and transmit electronic clinical
trial records. Reference FHIR \cite{hl7-fhir} and DICOM \cite{dicom} as
interoperability standards. Explain how Physical AI introduces new categories
of electronic records (robot telemetry, sensor data, autonomous decision logs)
that must comply with Part 11 requirements.]

\subsubsection{Privacy and Security}
[Cover HIPAA Privacy and Security Rules (45 CFR Parts 160 and 164)
\cite{hipaa} as they apply to Physical AI systems handling protected health
information. Reference Physical AI Federated Learning \cite{fl-paper} and
federated learning in healthcare \cite{federated-learning-healthcare} as
privacy-preserving approaches.]

\subsection{Robotics-Specific Regulatory Considerations}
\label{subsec:robotics-regulatory}

[Adapt the robotics-specific content from research\_b/01\_research\_b\_part1.txt.
Cover FDA's position on robotically-assisted surgical (RAS) systems in oncology,
noting that such systems have not received marketing authorization specifically
for cancer prevention/treatment. Explain the significance of oncology-related
endpoints (recurrence, disease-free survival, overall survival) for Physical AI
trials.]

[MAJOR THEME: The regulatory landscape for Physical AI oncology trials is
complex but navigable. By building on existing frameworks rather than creating
entirely new ones, the adapted standards in this platform (Adaption: ICH
Harmonised Guideline \cite{ich-adapt}, Physical AI Adaption: 21 CFR Part 50
\cite{cfr50-adapt}, Physical AI Adaption: 21 CFR Part 312 \cite{cfr312-adapt})
provide a credible and practical path forward.]

\subsection{Gaps and Opportunities}
\label{subsec:gaps-opportunities}

[Identify regulatory gaps and ambiguities from research\_b for Physical AI
implementations. Cover: (a) lack of specific Physical AI trial guidance from
FDA; (b) jurisdictional overlaps between California and federal requirements;
(c) classification challenges for multi-function Physical AI systems;
(d) workforce transition considerations. For each gap, explain how this
National Platform addresses it through its comprehensive body of adapted
standards and documentation.]

[Reference research\_b/02\_research\_b\_part2.txt for supporting URLs and
citations. Connect gaps to solutions provided in subsequent sections of this
paper, particularly Sections~\ref{sec:ich-e6r3}, \ref{sec:cfr50},
\ref{sec:cfr312}, and \ref{sec:implementation}.]
